\chapter{Overview}
\label{ch:00-overview}

These are summary notes from EEML 2026, the Eastern European Machine Learning Summer School, held in
Cetinje, Montenegro from 27 July to 1 August 2026. They are built from the slide decks distributed
during the week, photographs of the two lectures where no deck was given, four tutorial notebooks,
and two files of personal notes. They are not a transcript. Each lecture chapter states what the
lecture is taken to be arguing, gives the three to six ideas that carry that argument with the
mathematics where the mathematics is the point, and closes by locating the lecture among the
others.

Everything substantive traces back to the materials. Connective reasoning that does not come from a
lecture is confined to marked boxes like the one at the end of this chapter, so it is always visible
which claims belong to a speaker and which do not.

\section*{The week}

\begin{center}
\small
\begin{tabularx}{\textwidth}{@{}l l X l@{}}
\toprule
\textbf{Day} & \textbf{Chapter} & \textbf{Lecture and speaker} & \textbf{Source} \\
\midrule
Mon 27 Jul & \ref{ch:01-intro-deep-learning} & Introduction to Deep Learning, Sarath Chandar
  & photos \\
           & \ref{ch:02-representational-geometry} & Representational Geometry and Alignment,
  Phillip Isola & deck \\
           & \ref{ch:03-diffusion-models} & Diffusion Models, Sander Dieleman & deck \\
Tue 28 Jul & \ref{ch:04-intro-reinforcement-learning} & Introduction to Reinforcement Learning,
  Andreea Deac & slides \\
           & \ref{ch:05-statistics-estimation} & Statistics: Learning to Estimate,
  Kyunghyun Cho & deck \\
Wed 29 Jul & \ref{ch:06-continual-learning} & Continual Learning, Clare Lyle & deck \\
           & \ref{ch:07-to-be-figured-out} & To Be Figured Out, Razvan Pascanu & deck \\
Thu 30 Jul & \ref{ch:08-ml-for-mathematics} & ML for Mathematics, Matija Tapu\v{s}kovi\'{c}
  & deck \\
Fri 31 Jul & \ref{ch:09-spectral-analysis-dags} & Spectral Analysis of DAGs,
  Isidora Stankovi\'{c} & deck \\
           & \ref{ch:10-efficient-learning} & Computationally-Efficient Learning, Dan Alistarh
  & photos \\
Sat 1 Aug  & \ref{ch:11-geometric-deep-learning} & Geometric Deep Learning,
  Petar Veli\v{c}kovi\'{c} & deck \\
\bottomrule
\end{tabularx}
\end{center}

\noindent Four tutorial notebooks accompanied the week, on reinforcement learning, multimodal
learning, mechanistic interpretability and model quantization. They are not written up here,
since the notebooks are self-contained and are better worked through directly. Where a lecture
chapter has a matching tutorial, it says so.

\section*{Where to start}

On a return to this material after some time, the three synthesis chapters at the end are the point
of the exercise and the lecture chapters are their evidence. Chapter~\ref{ch:themes} identifies the
six ideas that recurred, including a genuine disagreement between Lyle and Pascanu about how serious
the i.i.d.\ assumption is. Chapter~\ref{ch:research-directions} separates problems the speakers
flagged as open from directions inferred here. Chapter~\ref{ch:revisit} is an ordered queue of
specific items to return to, with a reason for each.

\begin{aside}
This box is the mechanism for marking editorial additions, so here is one about the notes
themselves. The two lectures with no distributed deck, Chapters~\ref{ch:01-intro-deep-learning}
and~\ref{ch:10-efficient-learning}, are reconstructed entirely from photographs of projected slides,
and although the transcriptions were made at full resolution and are legible, those two chapters
rest on a single source with nothing to corroborate them. The lectures with both a deck and
photographs are the best evidenced, and there the photographs were left untranscribed because the
deck supersedes them.
\end{aside}
